% Part: first-order-logic
% Chapter: completeness
% Section: complete-sets

% Definition of complete consistent sets. Properties of complete sets required
% for completeness proved are in provability.tex in the
% chapter on the proof system used.

\documentclass[../../../include/open-logic-section]{subfiles}

\begin{document}

\iftag{FOL}
      {\olfileid{fol}{com}{ccs}}
      {\olfileid{pl}{com}{ccs}}

\olsection{\usetoken{P}{sentence}的完全一致集}

\begin{defn}[完全集]
\ollabel{def:complete-set} 语句集~$\Gamma$ 是\emph{完全的}，当且仅当对任意语句~$!A$，要么 $!A \in \Gamma$，要么 $\lnot !A \in \Gamma$。
\end{defn}

\begin{explain}
完全的语句集不留任何未决的问题。对于任何语句~$!A$，$\Gamma$ 都会“指出”$!A$ 为真还是为假。完全集的重要性不限于完备性定理的证明。例如，一个完全且可公理化的理论总是可判定的。
\end{explain}

\begin{explain}
完全一致集在完备性证明中非常重要，因为我们可以保证，每个一致的语句集~$\Gamma$ 都包含于某个完全一致集~$\Gamma^*$ 中。一个完全一致集对每个语句~$!A$ 都包含 $!A$ 或其否定 $\lnot !A$，但不同时包含两者。这对于\iftag{FOL}{原子语句}{命题变元}尤其成立，因此从一个完全一致集\iftag{FOL}{（在经适当添加常项符号而扩充的语言中）}{}，我们可以构造一个\iftag{FOL}{结构}{赋值}，其中\iftag{FOL}{谓词符号的解释}{指派给命题变元的真值}是根据哪些\iftag{FOL}{原子语句}{命题变元}在~$\Gamma^*$ 中来定义的。然后可以证明，这个\iftag{FOL}{结构}{赋值}使~$\Gamma^*$ 中的所有语句（从而也使~$\Gamma$ 中的所有语句）为真。证明后一个事实需要：$\lnot !A \in \Gamma^*$ 当且仅当 $!A \notin \Gamma^*$，$(!A \lor !B) \in
\Gamma^*$ 当且仅当 $!A \in \Gamma^*$ 或 $!B \in \Gamma^*$，等等。
\end{explain}

在下文中，我们将经常默认使用 $\Proves$ 的自反性、单调性和传递性（参见\iftag{FOL}{%
  \tagrefs{prfSC/{fol:seq:ptn:sec},prfND/{fol:ntd:ptn:sec},prfAX/{fol:axd:ptn:sec},prfTab/{fol:tab:ptn:sec}}}{%
  \tagrefs{prfSC/{pl:seq:ptn:sec},prfND/{pl:ntd:ptn:sec},prfAX/{pl:axd:ptn:sec},prfTab/{pl:tab:ptn:sec}}}）。

\begin{prop}
\ollabel{prop:ccs}
设 $\Gamma$ 是完全且一致的。则：
\begin{enumerate}
\item \ollabel{prop:ccs-prov-in} 若 $\Gamma \Proves !A$，则 $!A \in \Gamma$。

\tagitem{prvAnd}{\ollabel{prop:ccs-and} $!A \land !B \in \Gamma$ 当且仅当 $!A \in \Gamma$ 且 $!B \in \Gamma$。}{}

\tagitem{prvOr}{\ollabel{prop:ccs-or} $!A \lor !B \in \Gamma$ 当且仅当 $!A \in \Gamma$ 或 $!B \in \Gamma$。}{}

\tagitem{prvIf}{\ollabel{prop:ccs-if} $!A \lif !B \in \Gamma$ 当且仅当 $!A \notin \Gamma$ 或 $!B \in \Gamma$。}{}
\end{enumerate}
\end{prop}

\begin{proof}
以下我们始终假设 $\Gamma$ 是完全且一致的。
\begin{enumerate}
\item 若 $\Gamma \Proves !A$，则 $!A \in \Gamma$。

假设 $\Gamma \Proves !A$。为了反证，设 $!A
\notin \Gamma$。由于 $\Gamma$ 是完全的，故 $\lnot !A \in \Gamma$。
根据
\iftag{FOL}{%
  \tagrefs{prfSC/{fol:seq:prv:prop:explicit-inc},
    prfND/{fol:ntd:prv:prop:explicit-inc},
    prfAX/{fol:axd:prv:prop:explicit-inc},
    prfTab/{fol:tab:prv:prop:explicit-inc}}}{%
  \tagrefs{prfSC/{pl:seq:prv:prop:explicit-inc},
    prfND/{pl:ntd:prv:prop:explicit-inc},
    prfAX/{pl:axd:prv:prop:explicit-inc},
    prfTab/{pl:tab:prv:prop:explicit-inc}}}，
$\Gamma$ 是不一致的。这与 $\Gamma$ 一致的假设矛盾。因此，$!A \notin
\Gamma$ 不可能成立，所以 $!A \in \Gamma$。

\tagitem{defAnd}{}{%
\iftag{probAnd}{练习。}{%
$!A \land !B \in \Gamma$ 当且仅当 $!A \in \Gamma$ 且 $!B \in \Gamma$：

先证正向：设 $!A \land !B \in \Gamma$。则根据
\iftag{FOL}{%
  \tagrefs{prfSC/{fol:seq:ppr:prop:provability-land},
    prfND/{fol:ntd:ppr:prop:provability-land},
    prfAX/{fol:axd:ppr:prop:provability-land},
    prfTab/{fol:tab:ppr:prop:provability-land}}}{%
  \tagrefs{prfSC/{pl:seq:ppr:prop:provability-land},
    prfND/{pl:ntd:ppr:prop:provability-land},
    prfAX/{pl:axd:ppr:prop:provability-land},
    prfTab/{pl:tab:ppr:prop:provability-land}}} 第 (1) 条，
有 $\Gamma \Proves !A$ 且 $\Gamma \Proves !B$。由 \olref{prop:ccs-prov-in}，$!A \in \Gamma$ 且 $!B \in \Gamma$，如所需。

再证反向：设 $!A \in \Gamma$ 且 $!B \in
\Gamma$。根据
\iftag{FOL}{%
  \tagrefs{prfSC/{fol:seq:ppr:prop:provability-land},
    prfND/{fol:ntd:ppr:prop:provability-land},
    prfAX/{fol:axd:ppr:prop:provability-land},
    prfTab/{fol:tab:ppr:prop:provability-land}}}{%
  \tagrefs{prfSC/{pl:seq:ppr:prop:provability-land},
    prfND/{pl:ntd:ppr:prop:provability-land},
    prfAX/{pl:axd:ppr:prop:provability-land},
    prfTab/{pl:tab:ppr:prop:provability-land}}} 第 (2) 条，
有 $\Gamma \Proves !A \land !B$。再由 \olref{prop:ccs-prov-in}，得 $!A \land
!B \in \Gamma$。}}

\tagitem{defOr}{}{%
\iftag{probOr}{练习。}{%
先证若 $!A \lor !B \in \Gamma$，则 $!A \in
\Gamma$ 或 $!B \in \Gamma$。设 $!A \lor !B \in \Gamma$，但 $!A
\notin \Gamma$ 且 $!B \notin \Gamma$。由于 $\Gamma$ 是完全的，故 $\lnot !A \in \Gamma$ 且 $\lnot !B \in \Gamma$。根据
\iftag{FOL}{%
  \tagrefs{prfSC/{fol:seq:ppr:prop:provability-lor},
    prfND/{fol:ntd:ppr:prop:provability-lor},
    prfAX/{fol:axd:ppr:prop:provability-lor},
    prfTab/{fol:tab:ppr:prop:provability-lor}}}{%
  \tagrefs{prfSC/{pl:seq:ppr:prop:provability-lor},
    prfND/{pl:ntd:ppr:prop:provability-lor},
    prfAX/{pl:axd:ppr:prop:provability-lor},
    prfTab/{pl:tab:ppr:prop:provability-lor}}}
第 (1) 条，$\Gamma$ 不一致，矛盾。因此，$!A \in \Gamma$ 或 $!B \in \Gamma$。

再证反向：设 $!A \in \Gamma$ 或 $!B \in
\Gamma$。根据
\iftag{FOL}{%
  \tagrefs{prfSC/{fol:seq:ppr:prop:provability-lor},
    prfND/{fol:ntd:ppr:prop:provability-lor},
    prfAX/{fol:axd:ppr:prop:provability-lor},
    prfTab/{fol:tab:ppr:prop:provability-lor}}}{%
  \tagrefs{prfSC/{pl:seq:ppr:prop:provability-lor},
    prfND/{pl:ntd:ppr:prop:provability-lor},
    prfAX/{pl:axd:ppr:prop:provability-lor},
    prfTab/{pl:tab:ppr:prop:provability-lor}}} 第 (2) 条，
有 $\Gamma \Proves !A \lor !B$。由 \olref{prop:ccs-prov-in}，得 $!A \lor
!B \in \Gamma$，如所需。}}

\tagitem{defIf}{}{%
\iftag{probIf}{练习。}{%
先证正向：设 $!A \lif !B \in \Gamma$，并为了反证，设 $!A \in \Gamma$ 且 $!B \notin \Gamma$。在这些假设下，有 $!A \lif !B \in \Gamma$ 和 $!A \in \Gamma$。根据
\iftag{FOL}{%
  \tagrefs{prfSC/{fol:seq:ppr:prop:provability-lif},
    prfND/{fol:ntd:ppr:prop:provability-lif},
    prfAX/{fol:axd:ppr:prop:provability-lif},
    prfTab/{fol:tab:ppr:prop:provability-lif}}}{%
  \tagrefs{prfSC/{pl:seq:ppr:prop:provability-lif},
    prfND/{pl:ntd:ppr:prop:provability-lif},
    prfAX/{pl:axd:ppr:prop:provability-lif},
    prfTab/{pl:tab:ppr:prop:provability-lif}}} 第 (1) 条，
$\Gamma \Proves !B$。但由 \olref{prop:ccs-prov-in}，则有 $!B \in
\Gamma$，这与 $!B \notin \Gamma$ 的假设矛盾。

再证反向：首先考虑 $!A \notin
\Gamma$ 的情形。由于 $\Gamma$ 是完全的，故 $\lnot !A \in \Gamma$。根据
\iftag{FOL}{%
  \tagrefs{prfSC/{fol:seq:ppr:prop:provability-lif},
    prfND/{fol:ntd:ppr:prop:provability-lif},
    prfAX/{fol:axd:ppr:prop:provability-lif},
    prfTab/{fol:tab:ppr:prop:provability-lif}}}{%
  \tagrefs{prfSC/{pl:seq:ppr:prop:provability-lif},
    prfND/{pl:ntd:ppr:prop:provability-lif},
    prfAX/{pl:axd:ppr:prop:provability-lif},
    prfTab/{pl:tab:ppr:prop:provability-lif}}} 第 (2) 条，
$\Gamma \Proves !A \lif !B$。再次由 \olref{prop:ccs-prov-in}，得 $!A \lif !B \in \Gamma$，如所需。

现考虑 $!B \in \Gamma$ 的情形。根据
\iftag{FOL}{%
  \tagrefs{prfSC/{fol:seq:ppr:prop:provability-lif},
    prfND/{fol:ntd:ppr:prop:provability-lif},
    prfTab/{fol:tab:ppr:prop:provability-lif},
    prfAX/{fol:axd:ppr:prop:provability-lif}}}{%
  \tagrefs{prfSC/{pl:seq:ppr:prop:provability-lif},
    prfND/{pl:ntd:ppr:prop:provability-lif},
    prfTab/{pl:tab:ppr:prop:provability-lif},
    prfAX/{pl:axd:ppr:prop:provability-lif}}}
第 (2) 条同样有 $\Gamma \Proves !A \lif !B$。由
\olref{prop:ccs-prov-in}，得 $!A \lif !B \in \Gamma$。}}
\end{enumerate}
\end{proof}

\tagprob{FOL}


\begin{prob}
完成 \olref[fol][com][ccs]{prop:ccs} 的证明。
\end{prob}


\tagendprob

\tagprob{notFOL}


\begin{prob}
完成 \olref[pl][com][ccs]{prop:ccs} 的证明。
\end{prob}


\tagendprob

\end{document}